\chapter{KẾT QUẢ MÔ PHỎNG VÀ THẢO LUẬN}
\label{chap:results}

\section{Mục tiêu phân tích kết quả}
Chương này trả lời bốn câu hỏi nghiên cứu đã nêu trong phần Mở đầu. Trọng tâm không chỉ là phương pháp nào đạt giá trị trung bình lớn hơn, mà còn là: kết quả có đi kèm QoS hay không; chênh lệch có ổn định trên các scenario hay không; thời gian ra quyết định có phù hợp với mục tiêu suy luận nhanh hay không; và kết luận nào được phép suy rộng từ mô hình hiện tại.

Bộ kết quả cuối được audit trên năm kích thước $N\in\{16,32,64,96,128\}$. TD3 có tám seed độc lập cho mỗi $N$, mỗi seed được đánh giá trên 1.000 test scenario khóa. Các baseline xác định được chạy trên cùng 1.000 scenario. Coverage, giá trị hữu hạn, checksum, provenance, bảng, hình và benchmark CPU đã qua kiểm tra. Các con số trong chương được lấy từ bundle kết quả đã khóa; không sử dụng kết quả cũ của MADDPG hoặc bất kỳ phiên bản môi trường trước đó.

\section{Tiêu chí diễn giải}
Khi đọc kết quả, cần phân biệt bốn tầng:
\begin{enumerate}[label=\arabic*)]
  \item \textbf{Chất lượng phổ:} mean sum-rate và chênh lệch ghép cặp.
  \item \textbf{Độ tin cậy dịch vụ:} QoS fraction, all-users-QoS probability và violation.
  \item \textbf{Chi phí ra quyết định:} inference hoặc solve latency.
  \item \textbf{Độ chắc chắn thống kê:} CI, paired tests, effect size và giới hạn số seed.
\end{enumerate}
Không phương pháp nào được xem là tốt nhất tuyệt đối nếu chỉ thắng một tầng. AnalyticalRIS có thể nhanh nhất nhưng QoS yếu; AO--SCA có sum-rate cao nhất nhưng chậm; TD3 tạo điểm cân bằng.

\section{Tổng hợp kết quả chính}
Bảng~\ref{tab:audited-findings} tóm tắt các phát hiện đã được reviewer nội bộ chấp nhận cho tích hợp luận văn.

\begin{table}[H]
\caption{Các phát hiện định lượng chính từ bundle kết quả đã kiểm toán}
\label{tab:audited-findings}
\begin{tabularx}{\textwidth}{p{5.2cm}Y}
\toprule
\textbf{Chỉ số} & \textbf{Kết quả đã kiểm toán}\\
\midrule
Kích thước STAR--RIS & $N=16,32,64,96,128$\\
Số seed TD3 & 8 seed độc lập cho mỗi $N$\\
Số test scenario & 1.000 scenario khóa cho mỗi seed và mỗi $N$\\
Mean QoS fraction của TD3 & Từ 0,9892 đến 0,9994\\
All-users-QoS probability của TD3 & Từ 0,9798 đến 0,9985\\
Mean sum-rate TD3 & 10,7661 bit/s/Hz tại $N=16$; 12,4117 bit/s/Hz tại $N=128$\\
TD3 so với AO--Grid & Cao hơn 169,1\% đến 211,7\% về sum-rate\\
TD3 so với AnalyticalRIS & Cao hơn 440,7\% đến 526,2\% về sum-rate\\
AO--SCA so với TD3 & AO--SCA cao hơn 4,4476 đến 7,2246 bit/s/Hz\\
Latency TD3 so với AO--SCA & TD3 nhanh hơn 884 đến 3690 lần\\
Latency TD3 so với AO--Grid & TD3 nhanh hơn 299 đến 1748 lần\\
AnalyticalRIS so với TD3 & AnalyticalRIS nhanh hơn khoảng 2,07 đến 2,30 lần, nhưng không thỏa QoS trong thí nghiệm\\
Paired $t$-test + Holm & 15/15 so sánh có ý nghĩa ở $\alpha=0,05$\\
Exact Wilcoxon + Holm & 0/15 so sánh có ý nghĩa ở $\alpha=0,05$\\
\bottomrule
\end{tabularx}
\tablesource{Tác giả tổng hợp từ \texttt{results/FINAL\_RESULTS\_REVIEW.md} (2026)}
\end{table}

Bảng cho thấy thông điệp khoa học trung tâm: TD3 không thay thế AO--SCA về chất lượng sum-rate, nhưng cung cấp QoS rất cao và latency thấp hơn nhiều. Vì vậy, cách viết ``TD3 vượt mọi phương pháp'' là sai. Cách viết phù hợp là ``TD3 tạo trade-off chất lượng--độ trễ có lợi trong mô hình đã xét''.

\section{Phân tích tổng tốc độ}
\subsection{Xu hướng theo số phần tử STAR--RIS}
Mean sum-rate TD3 tăng từ 10,7661 tại $N=16$ lên 12,4117 tại $N=128$, tương đương mức tăng tuyệt đối 1,6456 bit/s/Hz. Xu hướng này phù hợp trực giác: số phần tử lớn hơn tạo thêm bậc tự do để định hình đường ghép tầng. Tuy nhiên, tốc độ tăng không nhất thiết tuyến tính theo $N$, vì cùng một STAR--RIS phải phục vụ nhiều người dùng và common rate bị giới hạn bởi user yếu nhất.

Không nên diễn giải mức tăng như ``array gain $N^2$''. Mô hình chứa kênh Gaussian không chuẩn hóa theo aperture, không có path loss theo kích thước bề mặt và không mô hình điện từ. Các cảnh báo trong tài liệu RIS cho thấy claim về scaling vật lý phải dựa trên mô hình diện tích và suy hao nhất quán \cite{Bjornson2020RISMyths,DiRenzo2020SmartRadio}. Ở đây, kết luận chỉ là trong generator và normalization đã khai báo, policy khai thác được số phần tử tăng để cải thiện metric mô phỏng.

\subsection{TD3 so với AO--Grid}
TD3 vượt AO--Grid từ 169,1\% đến 211,7\% về sum-rate. AO--Grid dùng codebook tọa độ hữu hạn; khi số chiều lớn, tìm kiếm theo từng tọa độ và độ phân giải thô có thể bỏ qua tương tác giữa nhiều pha. TD3 học ánh xạ toàn cục từ CSI sang toàn vector action, cho phép thay đổi phối hợp nhiều tọa độ trong một lần inference.

Kết quả không chứng minh mọi grid search đều yếu hơn TD3. Một codebook dày hơn hoặc chiến lược coordinate search khác có thể cải thiện chất lượng nhưng làm tăng số evaluation. Do đó AO--Grid trong luận văn là baseline xác định với codebook đã khai báo, không đại diện upper bound của họ phương pháp tìm kiếm rời rạc.

\subsection{TD3 so với AnalyticalRIS}
TD3 vượt AnalyticalRIS từ 440,7\% đến 526,2\%. Chênh lệch lớn có thể giải thích bởi AnalyticalRIS chỉ căn chỉnh pha tổng hợp và dùng phân bổ đều cho beta, công suất và common rate. Trong hệ nhiều người dùng có QoS, equal allocation hiếm khi tối ưu vì độ lợi kênh không đồng đều. Actor TD3 có thể tăng common power hoặc common fraction cho user yếu, đồng thời điều chỉnh beta theo hai phía.

Tuy nhiên, baseline quá đơn giản có thể làm tỷ lệ cải thiện phần trăm trở nên rất lớn. Với mục tiêu công bố Q1, cần tránh dùng mức 440--526\% như bằng chứng duy nhất về ưu thế của TD3. Baseline mạnh hơn là AO--SCA, và chính baseline này đạt sum-rate cao hơn TD3.

\subsection{TD3 so với AO--SCA}
AO--SCA vượt TD3 từ 4,4476 đến 7,2246 bit/s/Hz. Đây là chênh lệch đáng kể và phải được báo cáo nổi bật. AO--SCA đánh giá trực tiếp merit của từng scenario, ước lượng gradient theo biến vật lý và backtracking để bảo đảm không giảm. TD3 chỉ thực hiện một forward pass của mạng đã huấn luyện và không tinh chỉnh theo scenario, vì vậy tồn tại amortization gap.

Khoảng cách này có ba nguồn có thể có. Thứ nhất, actor MLP hữu hạn không biểu diễn được chính xác ánh xạ tối ưu phức tạp và có tính đa cực trị. Thứ hai, critic approximation và exploration có thể không khám phá được vùng action tốt. Thứ ba, reward ưu tiên QoS và checkpoint selection feasibility-first có thể hy sinh sum-rate để đạt độ tin cậy cao. Kết quả vì vậy không phải thất bại của TD3, mà phản ánh trade-off giữa tối ưu online và suy luận nhanh.

\section{Phân tích QoS}
\subsection{QoS fraction}
Mean QoS fraction TD3 nằm trong khoảng 0,9892--0,9994. Với $K=4$, metric mỗi scenario nhận các giá trị $0$, $0,25$, $0,5$, $0,75$, $1$. Trung bình gần một cho thấy đa số scenario có tất cả hoặc gần tất cả người dùng thỏa $R_{\min}=0,5$ bit/s/Hz.

Khoảng trên không nói rõ $N$ nào đạt cực tiểu/cực đại, do bảng reviewer chỉ báo range. Luận văn không tự gán giá trị cho từng $N$ khi raw table chưa được trích trực tiếp. Cách báo cáo range giữ tính trung thực và tránh nội suy.

\subsection{All-users-QoS probability}
All-users-QoS probability từ 0,9798 đến 0,9985. Đây là kết quả mạnh hơn QoS fraction vì yêu cầu toàn bộ bốn user cùng thỏa. Mức thấp nhất 0,9798 tương ứng khoảng 20 scenario trên 1.000 có ít nhất một user vi phạm, nếu xét một evaluation aggregate. Mức cao nhất 0,9985 tương ứng khoảng 1--2 scenario trên 1.000.

Sự khác biệt giữa QoS fraction và all-QoS cho thấy một số vi phạm tập trung ở ít user hoặc ít scenario. Việc báo cáo cả hai metric cần thiết: QoS fraction có thể che giấu trường hợp một user liên tục bị bỏ lại, còn all-QoS nhạy với bất kỳ vi phạm nào.

\subsection{Vai trò của dual controller và checkpoint selection}
QoS cao là kết quả của ba lớp bảo vệ: penalty tuyến tính và bình phương; multiplier dual thích nghi; và checkpoint selection feasibility-first. Nếu chỉ tối đa hóa sum-rate, policy có thể dành công suất cho user có kênh mạnh. Common rate allocation và dual penalty tạo động lực bù cho user yếu.

Tuy nhiên, không thể tách định lượng đóng góp của từng cơ chế chỉ từ bảng tổng hợp. Ablation riêng cho fixed penalty, không dual hoặc selection theo reward là cần thiết để kết luận nhân quả. Trong luận văn hiện tại, đây được xem là giải thích theo thiết kế source, không phải kết quả ablation đã hoàn tất.

\section{Phân tích độ trễ}
\subsection{TD3 và AO--SCA}
TD3 nhanh hơn AO--SCA từ 884 đến 3690 lần trong benchmark CPU một luồng. AO--SCA phải đánh giá finite-difference gradient cho từng tọa độ, lặp hai khối và có backtracking. Khi $N$ tăng, số biến STAR--RIS là $3N$, nên số evaluation tăng mạnh. TD3 chỉ thực hiện các phép nhân ma trận của actor và một lần tính metric.

Tỷ lệ speed-up không được chuyển thành claim về millisecond tuyệt đối trên mọi máy. Nó phụ thuộc CPU, BLAS, compiler, số thread, warm-up và implementation. Giá trị khoa học là so sánh trong một giao thức kiểm soát, nơi tất cả phương pháp chạy cùng chế độ một thread.

\subsection{TD3 và AO--Grid}
TD3 nhanh hơn AO--Grid từ 299 đến 1748 lần. AO--Grid không tính gradient nhưng phải đánh giá nhiều candidate theo codebook. Tăng độ mịn codebook sẽ tiếp tục tăng chi phí. Kết quả cho thấy lợi thế của amortized inference không chỉ so với SCA mà còn so với tìm kiếm xác định.

\subsection{TD3 và AnalyticalRIS}
AnalyticalRIS nhanh hơn TD3 khoảng 2,07--2,30 lần. Điều này hợp lý vì công thức căn chỉnh pha và equal allocation đơn giản hơn forward pass MLP. Do đó không được gọi TD3 là phương pháp nhanh nhất tuyệt đối. Lợi thế của TD3 là chất lượng và QoS cao hơn nhiều với latency vẫn thấp.

Đây là ví dụ điển hình của đường biên Pareto. AnalyticalRIS nằm ở cực latency thấp nhưng chất lượng thấp; AO--SCA nằm ở cực chất lượng cao nhưng latency lớn; TD3 nằm giữa. Trong ứng dụng cần quyết định rất nhanh và có QoS, TD3 có thể là lựa chọn hợp lý. Trong tối ưu offline không giới hạn thời gian, AO--SCA có thể được ưu tiên.

\section{Phân tích thống kê}
\subsection{Mâu thuẫn giữa paired $t$-test và Wilcoxon}
Paired $t$-test sau Holm cho 15/15 so sánh có ý nghĩa, nhưng exact Wilcoxon sau Holm cho 0/15. Hai kết quả khác nhau có thể xuất phát từ cấu trúc chênh lệch, ties, phân bố không đối xứng hoặc số lượng đơn vị độc lập sau aggregation. Paired $t$-test nhạy với mean difference, trong khi Wilcoxon sử dụng thứ hạng và giả định về đối xứng khác.

Do đó, câu ``kết quả có ý nghĩa thống kê theo cả hai kiểm định'' là không đúng. Luận văn báo cáo cả hai và dựa thêm vào effect size, mean difference và CI. Nếu mean difference lớn nhưng Wilcoxon--Holm không qua ngưỡng, bằng chứng về chênh lệch trung bình vẫn tồn tại nhưng độ bền phi tham số chưa đủ.

\subsection{Đơn vị độc lập và số seed}
Raw TD3 có tám seed trên mỗi scenario. Nếu xem mọi dòng seed--scenario là độc lập, sample size bị phóng đại vì các dòng chia sẻ cùng scenario. Pipeline trung bình seed theo method/scenario trước khi ghép cặp baseline, do đó đơn vị ghép cặp là scenario. Tuy nhiên, bất định do huấn luyện chỉ có tám seed và cần được thảo luận riêng.

Tám seed là tốt hơn một hoặc ba seed, nhưng vẫn hạn chế khi policy variance lớn. Một nghiên cứu mở rộng có thể dùng hierarchical bootstrap: lấy mẫu seed và scenario theo hai tầng, hoặc fit mixed-effects model để tách variance giữa seed và giữa scenario.

\subsection{Khoảng tin cậy cho metric bị chặn}
QoS fraction và all-QoS thuộc $[0,1]$. Student-$t$ CI có thể vượt biên khi mean gần 1. Báo cáo reviewer yêu cầu giữ raw values và dùng clipped display hoặc bounded/bootstrap interval. Trong luận văn, bảng chính dùng mean/range; khi đưa CI, nên ghi rõ phương pháp và không diễn giải phần vượt 1 như xác suất thực.

\section{Khả năng mở rộng}
\subsection{Mở rộng số chiều observation và action}
Khi $N$ tăng từ 16 lên 128, observation tăng từ 140 lên 1036 chiều và action từ 57 lên 393 chiều. Actor hidden width giữ 256, nên bottleneck biểu diễn trở nên mạnh hơn ở $N$ lớn. Dù vậy, TD3 vẫn duy trì QoS cao và sum-rate tăng, cho thấy blockwise normalization và dimension-aware noise có hiệu quả thực nghiệm.

\subsection{Giới hạn của kiến trúc MLP cố định}
MLP xử lý các phần tử STAR--RIS như một vector có thứ tự, không khai thác cấu trúc permutation hoặc locality. Khi $N$ vượt 128, số tham số đầu vào/đầu ra tăng tuyến tính và cần train riêng cho mỗi $N$. Kiến trúc set/graph/attention có thể chia sẻ tham số theo phần tử và hỗ trợ generalization giữa kích thước.

\subsection{Scaling của baseline}
AO--SCA chịu chi phí finite difference theo $d_a$ và mỗi evaluation theo $KN$, do đó độ trễ có xu hướng tăng nhanh hơn TD3. AO--Grid cũng tăng số candidate. AnalyticalRIS giữ chi phí gần $O(KN)$ và vẫn nhanh, nhưng không cải thiện allocation. Kết quả latency ratio tăng theo range lớn phù hợp với phân tích này.

\section{Ablation và diễn giải thành phần}
Repository định nghĩa NoRIS, FixedRIS, RandomRIS và Equal-Power. Về mặt giả thuyết:
\begin{itemize}
  \item NoRIS đo đóng góp của đường gián tiếp.
  \item FixedRIS đo giá trị của tối ưu beta/pha so với cấu hình trung tính.
  \item RandomRIS đo xem policy học được cấu trúc hay chỉ hưởng lợi từ thêm đường ngẫu nhiên.
  \item Equal-Power đo vai trò của power allocation, trong khi giữ RIS và common allocation học được.
\end{itemize}

Khi viết bản cuối sau khi có bảng ablation đầy đủ, nên báo cáo chênh lệch paired theo từng $N$, không chỉ một biểu đồ trung bình. Nếu Equal-Power giảm QoS mạnh, có thể kết luận power allocation là thành phần chính cho QoS; nếu FixedRIS giảm sum-rate, pha/beta là thành phần chính cho gain kênh. Hiện tại, chương giữ mô tả thiết kế và không tạo số liệu chưa có trong reviewer summary.

\section{Trả lời câu hỏi nghiên cứu}
\subsection{RQ1: TD3 có học được action khả thi và QoS ổn định?}
Có, trong phạm vi test bank đã khóa. Action decoder bảo đảm ràng buộc hình học, còn mean QoS fraction 0,9892--0,9994 và all-QoS 0,9798--0,9985 cho thấy policy duy trì QoS rất cao trên scenario chưa dùng để chọn checkpoint. Kết luận không mở rộng sang CSI sai, kênh geometry khác hoặc phần cứng thật.

\subsection{RQ2: TD3 so với baseline như thế nào?}
TD3 vượt rõ AO--Grid và AnalyticalRIS về sum-rate, nhưng thấp hơn AO--SCA 4,4476--7,2246 bit/s/Hz. Vì vậy TD3 không phải maximum-sum-rate method. Nó là learned approximation có chất lượng trung gian và QoS cao.

\subsection{RQ3: Khi $N$ tăng, hiệu năng thay đổi ra sao?}
Sum-rate TD3 tăng từ 10,7661 lên 12,4117 giữa hai đầu $N=16$ và $N=128$, trong khi QoS vẫn gần 1. Observation/action tăng lớn nhưng policy vẫn ổn định. Không đủ cơ sở để claim scaling law vật lý.

\subsection{RQ4: Trade-off chất lượng--độ trễ là gì?}
AO--SCA có chất lượng cao nhất nhưng TD3 nhanh hơn 884--3690 lần. AnalyticalRIS nhanh hơn TD3 khoảng hai lần nhưng không thỏa QoS. TD3 nằm ở vùng cân bằng: latency gần inference mạng nơ-ron, chất lượng thấp hơn solver mạnh nhưng cao hơn baseline đơn giản và có độ tin cậy QoS cao.

\section{Mối đe dọa đến tính hợp lệ}
\subsection{Internal validity}
Kết quả có thể bị ảnh hưởng bởi bug trong physics, merge hoặc benchmark. Rủi ro được giảm bằng test suite, shared physics, checksum, raw CSV và audit. Tuy nhiên, code audit không thay thế kiểm chứng độc lập hoặc formal verification.

\subsection{Construct validity}
Sum-rate Shannon và QoS threshold là proxy cho dịch vụ. Không có packet error, finite blocklength, fairness index, energy consumption hoặc signaling overhead. Latency đo solve time trên CPU, chưa bao gồm channel estimation và truyền CSI.

\subsection{External validity}
Kênh Gaussian tổng hợp và SISO giới hạn khả năng suy rộng. Các công trình trực tiếp thường dùng MISO, massive MIMO, bảo mật hoặc ISAC \cite{Ge2024RSMAMassiveMIMO,Chang2024CovertSTARRSMA,Liu2025STARRSMAISAC}. Pha độc lập cũng là giả định lý tưởng hóa \cite{Liu2021STARCoupledPhase}.

\subsection{Conclusion validity}
Mâu thuẫn giữa hai kiểm định và tám seed yêu cầu diễn giải thận trọng. Các tỷ lệ phần trăm lớn so với baseline yếu không đủ cho claim phổ quát. Kết luận chính dựa vào pattern nhất quán giữa quality, QoS và latency, không chỉ $p$-value.

% FIGURE-INTEGRATED: CH5 START
\begin{landscape}
\ThesisFigure{figures/chapter5_quality_latency_tradeoff.pdf}
{Không gian đánh đổi định tính giữa tổng tốc độ, độ tin cậy QoS và độ trễ ra quyết định}
{fig:ch5-tradeoff}[0.90\linewidth]
\figuresource{Tác giả xây dựng từ kết quả đã kiểm toán trong repository (2026)}
\end{landscape}

Hình~\ref{fig:ch5-tradeoff} tổng hợp định tính vị trí của bốn phương pháp trong không gian chất lượng--độ trễ. AnalyticalRIS nằm gần phía độ trễ thấp nhất nhưng được đánh dấu không đạt QoS; AO--Grid có chất lượng cao hơn cấu hình phân tích đơn giản nhưng vẫn thấp hơn TD3; TD3 nằm ở vùng trung gian với viền xanh biểu thị độ tin cậy QoS cao; AO--SCA nằm ở vùng tổng tốc độ cao nhất nhưng có thời gian giải lớn hơn. Mũi tên giữa TD3 và AO--SCA diễn giải hai chiều của trade-off: AO--SCA cải thiện chất lượng nghiệm, còn TD3 giảm mạnh độ trễ ra quyết định. Các tọa độ trong hình không phải số đo theo một thang tuyến tính và không thay thế bảng raw theo từng $N$; chúng chỉ tóm tắt kết luận đã được lượng hóa bằng các khoảng chênh lệch và tỷ lệ tốc độ ở các mục trước.
% FIGURE-INTEGRATED: CH5 END

\section{Hàm ý thực tiễn}
Trong một hệ thống có phân bố kênh tương đối ổn định và có thể huấn luyện offline, TD3 phù hợp làm bộ khởi tạo hoặc bộ ra quyết định nhanh. Khi thời gian cho phép, có thể dùng TD3 action làm khởi tạo cho AO--SCA, kỳ vọng giảm iteration và thu hẹp quality gap. Cấu trúc hybrid này kết hợp amortized inference với local refinement và là hướng nghiên cứu có giá trị.

Nếu distribution kênh thay đổi, cần cơ chế phát hiện drift và retraining. CSI không hoàn hảo có thể được xử lý bằng data augmentation hoặc robust observation. Với phần cứng phase-coupled, action decoder phải được thay bằng parameterization khả thi thay vì chỉ thêm penalty.

\section{Kết luận chương}
Kết quả cho thấy TD3 đạt QoS rất cao, cải thiện mạnh so với baseline đơn giản và có độ trễ thấp hơn nhiều so với solver lặp. Đồng thời, AO--SCA đạt sum-rate cao hơn, AnalyticalRIS nhanh hơn TD3 và hai kiểm định thống kê cho kết luận khác nhau. Những điểm này dẫn đến một kết luận cân bằng: TD3 là giải pháp low-latency, QoS-reliable với quality gap có thể đo được, không phải nghiệm tối ưu tuyệt đối. Đây là nền tảng cho phần Kết luận và hướng phát triển.
